\section{Middleware Composition}\label{sec:middleware}

Middleware in P-A-R addresses cross-cutting concerns that arise throughout the
agent runtime: structured logging for audit trails, automatic retry with
backoff for resilience, PII masking for compliance, and distributed tracing
for observability.  The central challenge is composing middleware so that (1)
interceptor execution order is predictable and composable, and (2) the
algebraic properties of the composition operator enable equational reasoning
about middleware pipelines.  We resolve this through the \emph{Russian Doll}
model, which treats each middleware as a wrapper around the call stack and
defines composition as a nest of layers.

\subsection{Hook Structure}\label{sec:hook-structure}

A \emph{middleware hook} is a record that optionally intercepts each stage of
agent execution.  Formally, let $\Ident$ be the set of conversational
contexts, $\call$ the set of tool-call records, and $\resp$ the set of LLM
responses.  A hook $h$ is a tuple

\[
h = \bigl(
  \onBeforeLLM,\;
  \onAfterLLM,\;
  \onBeforeTool,\;
  \onAfterTool,\;
  \onError
\bigr)
\]

where each field has type $\maybe{\alpha \to \maybe{\alpha}}$ as follows:

\[
\begin{array}{lcl}
\onBeforeLLM &\colon& \conv \to \conv\ \maybe \\
\onAfterLLM  &\colon& \resp \to \resp\ \maybe \\
\onBeforeTool &\colon& \call \to \call\ \maybe \\
\onAfterTool &\colon& \call \times \myresult \to \myresult\ \maybe \\
\onError &\colon& \myerror \to \myresult\ \maybe
\end{array}
\]

The special hook $\id$ (\emph{identity}) has all fields bound to $\None$.
Intercepting functions return $\Some{x}$ to continue the pipeline or
$\None$ to short-circuit with the (potentially modified) value.

\subsection{Russian Doll Composition}\label{sec:russian-doll}

The composition operator $\oplus$ layers two hooks so that $h_1$ wraps $h_2$.
Formally, for each interceptor field $f$ we define $(h_1 \oplus h_2).f$
pointwise using case analysis on whether either operand supplies a handler:

\[
(h_1 \oplus h_2).f =
\begin{cases}
\Some{\,x \mapsto h_2.f\,x\,} &\text{if } h_1.f = \None \land h_2.f \neq \None\\[4pt]
\Some{\,x \mapsto h_1.f\,x\,} &\text{if } h_1.f \neq \None \land h_2.f = \None\\[4pt]
\Some{\,x \mapsto h_1.f\,x >>= h_2.f\,} &\text{if } h_1.f \neq \None \land h_2.f \neq \None\\[4pt]
\None &\text{otherwise}
\end{cases}
\]

The monadic bind ($>>=$) sequences the two interceptors: the output of the
outer (earlier) handler is passed as input to the inner (later) handler.
Figure~\ref{fig:middleware-onion} illustrates the resulting concentric
structure.

\begin{figure}[t]
\centering
\begin{tikzpicture}[node distance=2mm, outer/.style={draw, rounded rectangle, minimum size=1.4cm, font=\ttfamily\small}]
% Concentric rings: outermost to innermost
\node[outer, fill=blue!20] (logging)    {\texttt{logging}};
\node[outer, fill=orange!20, above right=1.5cm and 1.5cm of logging] (retry) {\texttt{retry}};
\node[outer, fill=purple!20, above right=1.5cm and 1.5cm of retry] (pii)   {\texttt{pii\_mask}};
\node[outer, fill=green!20,  above right=1.5cm and 1.5cm of pii]   (core)  {\texttt{core}};

% Request flow arrow (left to right, outer to inner)
\draw[->, thick, blue] (-1.5cm, -0.6cm) -- node[above] {request flow} (logging.west);
\draw[->, thick, blue] (logging.east) .. controls +(+1.0cm, +0.3cm) .. (retry.west);
\draw[->, thick, blue] (retry.east)  .. controls +(+1.0cm, +0.3cm) .. (pii.west);
\draw[->, thick, blue] (pii.east)   .. controls +(+1.0cm, +0.3cm) .. (core.west);

% Response flow arrow (right to left, inner to outer)
\draw[->, thick, red] (core.east)   .. controls +(+1.0cm, -0.3cm) .. (pii.east);
\draw[->, thick, red] (pii.west)    .. controls +(-1.0cm, -0.3cm) .. (retry.east);
\draw[->, thick, red] (retry.west)  .. controls +(-1.0cm, -0.3cm) .. (logging.east);
\draw[->, thick, red] (1.5cm, 0.4cm) -- node[above] {response flow} (logging.east);

% Labels for interceptor directions
\node at (-2.2cm, -0.6cm) {\texttt{on\_before}};
\node at (2.2cm,  0.4cm) {\texttt{on\_after}};
\end{tikzpicture}
\caption{Russian Doll middleware composition.  Outer middleware executes
  \texttt{on\_before} first and \texttt{on\_after} last.}
\label{fig:middleware-onion}
\end{figure}

The OCaml implementation builds each interceptor by folding the hook list
from outermost to innermost (so the last hook in the list is the innermost
wrapper):

\begin{figure}[t]
\begin{lstlisting}[language=caml, numbers=left, stepnumber=1]
let apply_before_llm hooks conv next =
  fold_right (fun hook acc c ->
    match hook.on_before_llm with
    | Some f -> (match f c with Some c' -> acc c' | None -> acc c)
    | None -> acc c
  ) hooks next conv
\end{lstlisting}
\hfill
\caption{Composition of \texttt{on\_before\_llm} interceptors via
  \texttt{fold\_right}.  The last hook in the list executes closest to the
  core; the first hook in the list is the outermost wrapper.}
\label{fig:fold-right}
\end{figure}

Because \texttt{fold\_right} processes the hook list from right to left,
the \emph{first} middleware in the list becomes the \emph{outermost}
wrapper, executing its \texttt{on\_before} earliest and its
\texttt{on\_after} latest, which matches the Russian Doll intuition.

\subsection{Monoid Properties}\label{sec:monoid-properties}

The hook record with $\oplus$ forms a monoid.  We state all five properties
and sketch proofs.

\begin{description}
\item[Closure.]  For any hooks $h_1, h_2$, the record $h_1 \oplus h_2$ has
  all five fields defined (each case above produces a value in
  $\maybe{\alpha\to\maybe{\alpha}}$ or $\None$).  Hence $h_1 \oplus h_2$
  is a hook.

\item[Associativity.]  For all hooks $h_1, h_2, h_3$,
  $(h_1 \oplus h_2) \oplus h_3 = h_1 \oplus (h_2 \oplus h_3)$.
  Each field $f$ is defined pointwise using the monadic bind, which is
  associative: $(m_1 >>= f_1) >>= f_2 = m_1 >>= (f_1 >=> f_2)$.  Therefore the
  composition order is immaterial (validated empirically in
  \S\ref{sec:evaluation}).

\item[Identity.]  The hook $\id$ is the identity element:
  $h \oplus \id = \id \oplus h = h$ for all $h$.  When $h.f = \None$ the
  definition yields $\None$; when $h.f = \Some{g}$ the case
  $\id.f = \None$ and $h.f \neq \None$ yields $\Some{x \mapsto h.f\,x}$,
  which equals $h.f$ by $\beta$-reduction (validated empirically in
  \S\ref{sec:evaluation}).

\item[Interceptor Ordering -- \texttt{on\_before}.]  Let the hook list be
  $[h_1; h_2; \dots; h_n]$ where $h_1$ is the outermost.  The fold
  produces the nesting $((\dots(h_1 \circ h_2) \circ \dots) \circ h_n)$,
  so $h_1$'s interceptor runs \emph{first} on the way in.  The order of
  \texttt{on\_before} execution is $[h_1, h_2, \dots, h_n]$.

\item[Interceptor Ordering -- \texttt{on\_after}.]  By duality, the
  \texttt{on\_after} interceptors unwind from innermost to outermost, i.e.
  $[h_n, h_{n-1}, \dots, h_1]$.  This ordering is the reverse of
  \texttt{on\_before}, ensuring symmetric behaviour on the two legs of the
  call path.
\end{description}

These properties guarantee that middleware pipelines can be reordered,
split, and merged without altering behaviour, which is essential for
modular testing and safe composition of independently developed middleware.

\subsection{Example: Logging {\textrightarrow} Retry {\textrightarrow} PII Masking}\label{sec:middleware-example}

Consider the pipeline \texttt{[logging; retry; pii\_mask]} assembled in that
order (logging outermost, pii\_mask innermost).  Tracing a user request:

\begin{enumerate}
\item \texttt{logging.on\_before} emits a structured log entry recording the
  raw prompt and conversation ID, then passes the context to the next layer.
\item \texttt{retry.on\_before} annotates the context with retry metadata
  (attempt count, backoff state) and proceeds.
\item \texttt{pii\_mask.on\_before} scans the prompt for entities matching
  the PII regex; matching substrings are replaced with redacted tokens.
\item The core LLM call executes with the redacted, annotated context.
\item On receipt of the LLM response, \texttt{pii\_mask.on\_after} restores
  any redacted tokens to their original form using the stored masking map.
\item \texttt{retry.on\_after} checks the response for retryable status
  codes; if retry is required it re-invokes the pipeline, otherwise it
  passes the response outward.
\item \texttt{logging.on\_after} records the final response and elapsed
  time, then the pipeline terminates.
\end{enumerate}

Figure~\ref{fig:middleware-onion} depicts the concentric wrapping:
request flow (left arrows) penetrates from \texttt{logging} through
\texttt{retry} and \texttt{pii\_mask} into the core; response flow (right
arrows) unwinds in the reverse order.  The monoid properties from
\S\ref{sec:monoid-properties} guarantee that reordering the middleware
list, e.g.\ \texttt{[retry; logging; pii\_mask]}, yields a semantically
equivalent pipeline with appropriately shifted interceptor ordering.